\section{Harmonic forms and the Coefficient-Length conjecture}
\label{sec:length}

\subsection{The Euler splitting of harmonic forms}

A row hook has an alternating block of one-box rows below its first
row.  Replacing that block by exterior variables turns its
highest-weight condition into contraction by the Euler vector
field.  The de Rham identity then separates precisely the primitive
forms corresponding to row hooks.

On $K[x_1,\ldots,x_n]\otimes\Lambda(\theta_1,\ldots,\theta_n)$,
put
\[
 D_r=\sum_i\partial_{x_i}^r,\qquad
 \kappa_r=\sum_i\partial_{x_i}^r\partial_{\theta_i},
 \qquad
 \iota=\sum_i x_i\partial_{\theta_i},\qquad
 \dd=\sum_i\theta_i\partial_{x_i}.
\]
The superharmonic space $\SH_n$ is the joint kernel of $D_r$
for $r\geq1$ and $\kappa_r$ for $r\geq0$.
These are the apolar equations for the coinvariant algebra
of differential forms.  The invariant differential-form
description goes back to Solomon \cite{Solomon}.
Write $\SH_n^{a,j}$ for polynomial degree $a$ and exterior
degree $j$.

For $a\geq1$, let $K_{a,j}$ be the $S_n$-module of row
highest-weight vectors of weight $(a,1^j)$ in
$\Harm_{n,j+1}$.  For the boundary terms put
$K_{0,j}=0$ for $j>0$ and $K_{0,0}=K$.

\begin{lemma}[The row-hook bridge]\label{lem:hook-bridge}
There are $S_n$-equivariant identifications
\begin{equation}\label{eq:hook-bridge}
 K_{a,j}=\ker\iota\cap\SH_n^{a,j}.
\end{equation}
Away from bidegree $(0,0)$, there is a direct sum
\begin{equation}\label{eq:euler-splitting}
 \SH_n^{a,j}=K_{a,j}\oplus\dd K_{a+1,j-1},
\end{equation}
with nonexistent exterior degrees interpreted as zero.
\end{lemma}

\begin{proof}
Replace an exterior monomial by a determinant in $j$ auxiliary
rows:
\begin{equation}\label{eq:determinant-replacement}
 \theta_{i_1}\cdots\theta_{i_j}
 \longmapsto
 \det(y_{i_s}^{(r)})_{1\leq r,s\leq j}.
\end{equation}
This identifies exterior degree $j$ with polynomials
multilinear and alternating in the auxiliary rows.
Raising operators between auxiliary rows vanish on these
polynomials.  Raising an auxiliary row into the $x$-row
is precisely contraction by $\iota$.
Conversely, the auxiliary highest-weight conditions in
weight $(1,\ldots,1)$ select the determinant representation,
so every row-hook highest vector has this form.

The harmonic equations using only the $x$-row are $D_r$.
Those using exactly one auxiliary row are $\kappa_r$.
Any equation with two auxiliary derivatives at the same
label vanishes on an alternating determinant, and an equation
differentiating twice in an auxiliary row vanishes by
multilinearity.  These are all possible equations, proving
\eqref{eq:hook-bridge}.

Both $\iota$ and $\dd$ preserve $\SH_n$, since
\[
 [D_r,\dd]=0,\quad [D_r,\iota]=r\kappa_{r-1},\quad
 \kappa_r\dd+\dd\kappa_r=D_{r+1},\quad
 \kappa_r\iota+\iota\kappa_r=0.
\]
Moreover $\iota^2=\dd^2=0$ and
\[
 \iota\dd+\dd\iota=(a+j)\operatorname{id}
       \quad\hbox{on bidegree }(a,j).
\]
Thus every $f\in\SH_n^{a,j}$ with $a+j>0$ has the decomposition
\[
 f=\frac{\iota\dd f}{a+j}
       +\dd\left(\frac{\iota f}{a+j}\right).
\]
The first term belongs to $K_{a,j}$, and the expression in
parentheses belongs to $K_{a+1,j-1}$.  On this latter primitive
space, $\iota\dd=(a+j)\operatorname{id}$, so $\dd$ is injective
with inverse $\iota/(a+j)$ on its image.  The same identity shows
that its image meets $\ker\iota$ trivially.  This proves
\eqref{eq:euler-splitting}.
When $a=0$ and $j>0$, $\dd$ vanishes on the source;
the same identity gives $\ker\iota=0$, as required by
the boundary convention.
\end{proof}

\subsection{A tableau formula and generation of all row hooks}

For a standard tableau $T$, a descent is an integer $i$ for
which $i+1$ lies in a lower row than $i$.
Let $\des(T)$ be the number of descents and $\maj(T)$ their sum.

\begin{theorem}[Row-hook multiplicities]\label{thm:row-hooks}
For every $\mu\vdash n$ and $j\geq0$,
\begin{equation}\label{eq:row-hook-tableaux}
 \sum_{a\geq1}[K_{a,j}:S^\mu]q^a
 =
 \sum_{\substack{T\in\SYT(\mu)\\t=\des(T)\geq1}}
 q^{\maj(T)}e_j(q^{-2},q^{-3},\ldots,q^{-t}).
\end{equation}
Equivalently, the coefficient indexed by $(a,1^j)$ and $\mu$
counts pairs $(T,J)$ satisfying
\[
 J\subseteq\{2,\ldots,\des(T)\},\qquad
 |J|=j,\qquad \maj(T)-\sum_{r\in J}r=a.
\]
\end{theorem}

\begin{proof}
The Fields theorem of Murai, Rhoades and Wilson
\cite[Theorem~5.2 and equation~(2.5)]{MRW} gives
\begin{equation}\label{eq:fields-character}
 \Frob(\SH_n;q,u)
 =\sum_{\mu\vdash n}s_\mu[Z]
   \sum_{T\in\SYT(\mu)}
      q^{\maj(T)}
      \prod_{r=1}^{\des(T)}(1+u q^{-r}).
\end{equation}
The generating series of \eqref{eq:euler-splitting} says
that the nonconstant character is $(1+u/q)$ times the
generating series of the $K_{a,j}$.
The unique tableau without descents is the one-row tableau,
and contributes only the constant trivial representation.
For every other tableau the factor $1+u/q$ cancels.
Extracting $u^j$ gives \eqref{eq:row-hook-tableaux}.
All remaining powers of $q$ are positive:
if $t=\des(T)$, then
$\maj(T)\geq t(t+1)/2$, whereas
$\sum_{r=2}^t r=t(t+1)/2-1$.
\end{proof}

The harmonic-form bridge extracts the row-hook multiplicities from
the Fields character.  For the alternating label, the resulting
specialization also
appears in \cite[Corollary~4.6]{Lentfer}.
The row-hook coefficients after summing symmetric-group
dimensions are studied in \cite{CL}.

\begin{proposition}[Generation of the row-hook sector]
\label{prop:row-hook-generation}
Every row-hook isotypic component of $\Harm_{n,k}$ is
contained in $\Pol_{n,k}$.
\end{proposition}

\begin{proof}
For $r\geq1$ set $\Theta_r=\sum_i\theta_i\partial_{x_i}^r$.
The operator theorem of Rhoades and Wilson
\cite[Theorem~5.1]{RW} gives
\begin{equation}\label{eq:rw-operators}
 \SH_n=\sum_{J\subseteq[n-1]}\Theta_J\Harm_{n,1},
\end{equation}
where the $\Theta_r$ anticommute and a fixed order is used
in each product.
Every ordinary harmonic is a linear combination of partial
derivatives of $\Delta_n$.
Under \eqref{eq:determinant-replacement}, the form
$\Theta_{r_1}\cdots\Theta_{r_j}h$ becomes
\begin{equation}\label{eq:theta-polarization}
 \sum_{\sigma\in S_j}\sgn(\sigma)
 E_{r_1}^{(\sigma(1))}\cdots E_{r_j}^{(\sigma(j))}h,
 \qquad
 E_s^{(a)}=\sum_i y_i^{(a)}\partial_{x_i}^s.
\end{equation}
The even operators in this formula commute.
Expansion of both sides verifies the formula term by term.
It belongs to $\Pol_{n,j+1}$.
Thus the entire subspace
$\ker\iota\cap\SH_n^{a,j}$ is represented inside
$\Pol_{n,j+1}$.  By Lemma~\ref{lem:hook-bridge}, this
is the full row-hook highest-weight multiplicity space.
Closing under $\GL_{j+1}$, and then adjoining unused rows
and closing under $\GL_k$, proves the proposition.
\end{proof}

For a nontrivial hook label $\mu=(a+1,1^b)$, $a+b=n-1$, $b\ge1$,
the tableau formula and row-hook generation give
\begin{equation}\label{eq:hook-hook-formula}
 \frac{C_\mu^{\Pol}[q\mid u]}{q+u}
  =\begin{bmatrix}n-1\\b\end{bmatrix}_q
                 \prod_{r=2}^{b}(q^r+u).
\end{equation}
Here the supersymmetric specialization is characterized by
$p_r\mapsto q^r+(-1)^{r-1}u^r$, and the bracket is the Gaussian
binomial coefficient.  To verify the formula, every tableau of
this hook has $b$ descents, and their possible sets are all
$b$-element subsets of $[n-1]$.  Thus their major-index generating
polynomial is $q^{b(b+1)/2}\bigl[\begin{smallmatrix}n-1\\b\end{smallmatrix}\bigr]_q$.
Now use $s_{(c,1^j)}[q\mid u]=(q+u)q^{c-1}u^j$ and
Theorem~\ref{thm:row-hooks}.  All other nonempty row shapes vanish
under this specialization.  This proves the nontrivial hook-label
cases of Bergeron's Conjecture~4
\cite[Section~4.2, equation~(4.4)]{Bergeron}.

\subsection{Maximal row hooks and exact coefficient lengths}

At the longest possible hook length, the tableau formula forces the
maximum descent number. Transposing the tableaux turns this into a
minimum-descent condition, which is counted by a finite principal
specialization. This gives the whole multiplicity series at once.

\begin{theorem}[The complete maximal-length row-hook series]
\label{thm:maximal-hook-series}
Let $\mu=(c,\tau)\vdash n$ and $r=|\tau|>0$. Then, in both the
harmonic and Vandermonde modules,
\[
 H_\mu(q):=\sum_{a\ge1}[K_{a,r-1}:S^\mu]q^a
                         =q\,s_{\tau'}(1,q,\ldots,q^{c-1}).
\]
Thus all maximal-length row-hook coefficients are given by one finite
principal specialization of a Schur polynomial.
\end{theorem}
\begin{proof}
A tableau of shape $\mu$ has at most $r$ descents. In \eqref{eq:row-hook-tableaux}, exterior degree $r-1$ therefore forces exactly
$r$ descents and selects the whole set $\{2,\ldots,r\}$.
Transposition complements the descent set. If $N=\binom n2$, then
\[
 H_\mu(q)=q^{N-r(r+1)/2+1}
 \sum_{\substack{U\in\operatorname{SYT}(\mu')\\
                         \operatorname{des}(U)=c-1}}
      q^{-\operatorname{maj}(U)}
 =q^{N-r(r+1)/2+1}s_{\mu'}(1,q^{-1},\ldots,q^{-(c-1)}).
\]
For the second equality, use the standard principal-specialization
identity on exactly $c=\ell(\mu')$ variables. Standardization gives a
sum of $q^{\operatorname{maj}(U)}$ times Gaussian coefficients; only
the minimum descent number $c-1$ survives, with Gaussian coefficient
one. One may standardize using comajor index and use tableau evacuation
to replace comajor by major index.

Now $\mu'=(1^c)+\tau'$, with $\tau'$ padded by zeros. Factoring the
determinant contributes $q^{-\binom c2}$. Reverse the remaining $c$
variables and use homogeneity of $s_{\tau'}$, of degree $r$, to obtain
an additional factor $q^{-r(c-1)}$. Since
\[
 N-\frac{r(r+1)}2-\binom c2-r(c-1)=0,
\]
the total remaining factor is $q$, proving the formula.
\end{proof}

\begin{corollary}[A positive contribution at every permitted hook length]
\label{cor:shorter-hook-witnesses}
With $\mu=(c,\tau)$ and $r=|\tau|>0$, fix $1\le\ell\le r$.
In the row-hook coefficient series of length $\ell$, the contribution
of tableaux with exactly $r$ descents is
\[
 q^{\binom{r-\ell+2}{2}}
 \begin{bmatrix}r-1\\\ell-1\end{bmatrix}_q
                       s_{\tau'}(1,q,\ldots,q^{c-1}).
\]
All other contributions have nonnegative coefficients. Hence, if
$b=\sum_i\binom{\tau_i}{2}$, the row hook
\[
       \left(\binom{r-\ell+2}{2}+b,\,1^{\ell-1}\right)
\]
occurs. Taking $\ell=\min(k,r)$ gives the sharp coefficient length
in $k$ alphabets directly.
\end{corollary}
\begin{proof}
The proof of Theorem~\ref{thm:maximal-hook-series} gives
$\sum_{\operatorname{des}(T)=r}q^{\operatorname{maj}(T)}
=q^{r(r+1)/2}s_{\tau'}(1,q,\ldots,q^{c-1})$.
Multiply by
\[
 e_{\ell-1}(q^{-2},\ldots,q^{-r})
 =q^{-r(\ell-1)+\binom{\ell-1}{2}}
             \begin{bmatrix}r-1\\\ell-1\end{bmatrix}_q
\]
and simplify the exponent. Its least coefficient is positive.
\end{proof}




\begin{theorem}[Exact coefficient lengths]\label{thm:length}
For $\mu\vdash n$, put $r=n-\mu_1$.
The largest length of a row partition in the $\mu$-isotype of either
$\Harm_{n,k}$ or $\Pol_{n,k}$ is $\min(k,r)$.
The corresponding stable largest length is $r$.
\end{theorem}

\begin{proof}
Proposition~\ref{prop:length-upper} gives the upper bound. If
$\mu=(c,\tau)$ and $r=|\tau|>0$, take
$\ell=\min(k,r)$ in Corollary~\ref{cor:shorter-hook-witnesses}.
Its positive contribution contains a row hook of length $\ell$;
Proposition~\ref{prop:row-hook-generation} places the whole component
in $\Pol_{n,k}$. If $r=0$, exactness of taking $S_n$-invariants
makes the trivial isotype of $R_{n,k}$ equal to $K$, and constants
belong to $\Pol_{n,k}$.
\end{proof}

\begin{corollary}[Palindromicity, unimodality, and the extremal hook]
\label{cor:maximal-hook-coefficients}
Put $b=\sum_i\binom{\tau_i}{2}$. The coefficients of $H_\mu(q)$
are palindromic and unimodal on the interval
\[
             a_{\min}=1+b,\qquad a_{\max}=1+r(c-1)-b,
\]
with both endpoint coefficients equal to one. Their center is
$1+r(c-1)/2$. In particular the least-degree row hook of maximal
possible length $r$ is
\begin{equation}\label{eq:explicit-length-witness}
 \lambda_\mu=(1+b,1^{r-1}),\qquad c_{\lambda_\mu,\mu}^{(n)}=1.
\end{equation}
Its total degree is
$r+b=\sum_i\binom{\tau_i+1}{2}$. The sum of all maximal-length
row-hook multiplicities is
\[
 H_\mu(1)=\dim\mathbb S_{\tau'}K^c
 =\prod_{1\le i<j\le c}
                 \frac{\tau'_i-\tau'_j+j-i}{j-i}.
\]
An exact product form is
\[
 H_\mu(q)=q^{1+b}
   \prod_{(i,j)\in\tau}\frac{1-q^{c+i-j}}{1-q^{h_\tau(i,j)}}.
\]
\end{corollary}
\begin{proof}
The least principal-specialization exponent is
$n(\tau')=\sum_i\binom{\tau_i}{2}=b$, with coefficient one;
reversal of the $c$ variables gives the greatest exponent
$r(c-1)-b$, also with coefficient one. This proves the endpoints
and palindromicity.

For unimodality, apply the standard $\mathfrak{sl}_2$ criterion
\cite{Proctor1982,Stanley1989}: regard $K^c$ as the irreducible
$\mathfrak{sl}_2$-module $\operatorname{Sym}^{c-1}K^2$. Then
\[
 t^{-r(c-1)}s_{\tau'}(1,t^2,\ldots,t^{2(c-1)})
\]
is the $\mathfrak{sl}_2$ character of $\mathbb S_{\tau'}K^c$.
Its decomposition into irreducibles expresses its coefficients as
a sum of indicator functions of centered weight intervals. The
coefficients therefore increase weakly up to the center and decrease
weakly thereafter. Section~\ref{sec:hook-action} constructs the
operators on harmonic multiplicities. The Weyl dimension and hook-content
formulas give the last two identities.
\end{proof}

For $\mu=(3,3)$ this hook is $(4,1,1)$; for the alternating label
it is $(1^{n-1})$.

\begin{corollary}[Small row support and operator recovery]
\label{cor:near-row}
The stable character $C_\mu^{(n)}$ is determined by its
specialization to $n-\mu_1$ row variables.
If $\mu_1\geq n-2$, then the $\mu$-isotypes of
$\Pol_{n,k}$ and $\Harm_{n,k}$ agree for every $k$.
\end{corollary}

\begin{proof}
Schur functions of length at most $r$ remain linearly
independent in $r$ variables.  The length bound therefore
proves the first assertion.
For the second, Haiman's bivariate operator theorem
\cite[Theorem~4.2]{Haiman} gives $\Pol_{n,2}=\Harm_{n,2}$.
Every highest-weight vector in the specified labels
uses at most two rows, by the length bound and
Lemma~\ref{lem:stability}.  It is therefore generated
in those rows.  The row group generates its full
isotypic component in any larger number of rows.
The one-row assertion is classical.
\end{proof}

\begin{corollary}\label{cor:three-row-defect}
In three alphabets, a nonzero component of
$\Harm_{n,3}/\Pol_{n,3}$ can occur only at a pair
$(\lambda,\mu)$ satisfying
\[
 \lambda=(a,b,c),\qquad a\geq b\geq2,\quad c\geq1,
 \qquad \mu_1\leq n-3.
\]
\end{corollary}

\begin{proof}
All row partitions of length at most two are generated
by the bivariate theorem and trailing-zero stability.
The remaining row hooks are generated by
Proposition~\ref{prop:row-hook-generation}.
Corollary~\ref{cor:near-row} excludes all label types
with $\mu_1\geq n-2$.
\end{proof}

